\subsubsection{Geometrische Repräsentation und Rendering der Punktwolke}

Für die Visualisierung wurde eine punktwolkenbasierte Darstellung gewählt, da sie algorithmische Effizienz mit hoher visueller Detailtreue verbindet. Da die erfassten Geometriedaten bereits vorliegen, entfallen komplexe und rechenintensive Transformationen, wie sie bei Voxel- und Mesh-Generierung notwendig sind. Um den Rechenaufwand bei der Verarbeitung eines vollständigen Stadtgebiets zu begrenzen, beschränkt sich dieser Ansatz im Rahmen einer prototypischen Umsetzung auf die Visualisierung ausgewählter Abschnitte. Dies minimiert den Vorverarbeitungs- und Implementierungsaufwand erheblich. Zugleich garantiert dieser Ansatz eine hohe visuelle Authentizität, da morphologische Details ohne geometrische Abstraktion erhalten bleiben \cite{wangAutomatedPhenotypicTrait2024}. Das Rendering erfolgt dabei durch die \textit{PyVista}-Bibliothek.

Die räumliche Kontextualisierung erfolgt über ein DOP, das als Basistextur auf eine flache Ebene knapp unterhalb der Geländehöhe projiziert wird. Zur Unterstützung der visuellen Orientierung werden die extrahierten 2D-Kronenpolygone als Linien-Mesh über dem DOP eingeblendet. Eine dynamische Farbskala dient der visuellen Einordnung der attributbasierten Baumeinfärbung \cite{kempfObliqueViewIndividual2021}. Da standardmäßige, flache Pixel-Primitives keine Oberflächennormale besitzen und somit keine Beleuchtungsberechnung erlauben, wirken unoptimierte Punktwolken oft flach und strukturarm. Zur Steigerung der räumlichen Tiefe werden die LiDAR-Punkte als 3D-Spähren gerendert, die die Anwendung lokaler Beleuchtungsmodelle ermöglichen. Ergänzend wird ein Tiefenschatten mittels \textit{Eye-Dome-Lighting} integriert, welcher die Konturen der Baumsegmente und die relative Tiefenwahrnehmung zwischen den Punkten erheblich verstärkt \cite{wegenNonPhotorealisticRendering3D2024}.

Die Darstellung massiver Punktmengen steht jedoch im Konflikt mit einer flüssigen, interaktiven Exploration. Daher wird ein distanzabhängiges, fünfstufiges Level-of-Detail (\ac{LoD}) Verfahren implementiert \cite{schutzRealTimeContinuousLevel2019}, welches in Code \ref{lst:visualize_lod} dargestellt ist.
\begin{lstlisting}[label={lst:visualize_lod}, caption={[LoD-Visualisierung]Dynamische LoD-Steuerung der Punktwolke}, breaklines=true, breakatwhitespace=true]
def on_camera_changed(self, obj, event):
    #[...]
    camera = self.plotter.camera
    dist = camera.distance
    #[...]
    if not self.point_actor:
        return

    if dist < 250:
        new_tier = 0; target_size = 12; target_spheres = True; use_low_poly = False
    elif dist < 450:
        new_tier = 1; target_size = 10; target_spheres = True; use_low_poly = False
    elif dist < 700:
        new_tier = 2; target_size = 8; target_spheres = True; use_low_poly = False
    elif dist < 1200:
        new_tier = 3; target_size = 5; target_spheres = False; use_low_poly = True
    else:
        new_tier = 4; target_size = 3; target_spheres = False; use_low_poly = True

    if self.current_lod_tier != new_tier:
        self._block_lod = True
        try:
            self.current_lod_tier = new_tier
            dataset_to_use = self.cloud_low if use_low_poly else self.cloud
            #[...]        
            self.point_actor.mapper.dataset = dataset_to_use
            #[...]
            self.point_actor.GetProperty().SetRenderPointsAsSpheres (target_spheres)
            self.point_actor.GetProperty().SetPointSize(target_size)
            self.plotter.render()
        finally:
            self._block_lod = False
\end{lstlisting}
Der Algorithmus überwacht kontinuierlich die Kameradistanz (\textit{dist}) zum Fokuspunkt, um den Sichtraum in fünf Detailstufen (\textit{new\_tier}) zu unterteilen. Die Stufen \textit{0 bis 2} decken den Nah- und Mittelansicht bei einem Kameraabstand von weniger als \textit{700 m} ab. Hier wird die Punktwolke in maximaler Auflösung (\textit{self.cloud}) als geladen, die Punktegrößen \textit{(target\_size)} stufenweise verringert und der Parameter \textit{target\_spheres} auf \textit{True} gesetzt, um eine korrekte Belechtungsberechnung auf 3D-Sphären zu erzwingen. Bei größeren Abständen schaltet das System in den Fernbereich (\textit{Tier 3 und 4}) um. Durch das Setzten der der Flag \textit{use\_low\_poly} wird der Datensatz im Mapper dynamisch gegen eine mittels Downsampling ausgedünnte Punktwolke (\textit{self.cloud\_low}) ausgetauscht und als flache 2D-Primitives (\textit{target\_spheres = False}) gerendert. Diese daten- und geometrieseitige Reduktion minimiert die GPU-Last entscheidend und gewährleistet eine performante Visualisierung großflächiger Abschnitte. Um ein Blockieren der grafischen Benutzeroberfläche (\ac{GUI}) beim Einlesen umfangreicher Datenstrukturen zu verhindern, basiert die Architektur auf einem asynchronen Design, bei dem sämtliche Ladevorgänge in einen Hintergrund-Thread ausgelagert sind \cite{schutzSimLODSimultaneousLOD2024}.
